%قبل از شروع مسئله ترجیحاً چیزی ننویسید. اما اگر خواستید بنویسید، در قالب لیستی (فقرات و شمارش) بنویسید.
%برای اضافه کردن نمره، بعدی از شروع‌مسئله، [x] را اضافه کنید و به جای x  نمره مورد نظر را قرار دهید. برای زیرمسئله‌ها هم به همین صورت است.
%\شروع‌مسئله[10]

\شروع‌مسئله[10]
فرض کنید یک گراف بدون جهت داریم که هر یال آن دارای دو وزن مثبت است. بار اولی که از یک یال عبور می‌کنیم باید به اندازه ‌ی وزن بیشتر آن یال هزینه پرداخت کنیم و بارهای بعدی به اندازه‌ی وزن سبکتر هزینه پرداخت می‌کنیم. می‌خواهیم از راس u به راس v برویم و در مسیر از راس w عبور کنیم. الگوریتمی از مرتبه‌ی $O(n\log(n) + m)$ ارائه دهید که مسیر با کمترین وزن را پیدا کند که $n$ و $m$ تعداد رئوس و یال‌های گراف هستند.


%اگر مسئله چند قسمت داشت به این صورت عمل کنید:
%\شروع{زیرمسائل}
%\زیرمسئله[0] % مسئلۀ ۱آ
%\زیرمسئله % مسئلۀ ۱ب
%\زیرمسئله % مسئلۀ ۱پ
%\زیرمسئله % مسئلۀ ۱ت
%\پایان{زیرمسائل}

\solution{
ادعا می‌کنیم جواب بهینه حتما به شکل زیر است:
\\
\درج‌تصویر[پهنا=0.4\پهنای‌متن]{A1.PNG}
\\در جواب بهینه مسیری که از $u$ به $w$ طی کرده‌ایم را $p$ می‌نامیم. ادعا می‌کنیم بعد از رسیدن به $w$ یخشی از مسیر $p$ را باز‌میگردیم و سپس بدون برگشت به رئوس مسیر $p$ به راس $v$ می‌رویم. 
\\حال برای اثبات این ادعا، ثابت می‌کنیم بعد از عبور از راس $w$ اولین جایی که از مسیر $p$ خارج شویم دیگر به آن بر نمیگردیم. اولا اگر در مسیر $p$ از راسی دوبار عبور نمی‌کنیم. در مسیر از $w$ به $v$ نیز این شرط برقرار است چراکه در غیر این صورت در مسیر یک دور داریم که می‌توان آن دور را بدون مشکل از مسیر حذف کرد و مسیر کوتاه‌تر می‌شود (یعنی هر مسیر راس تکراری در بین رئوس خود آن مسیر ندارد). حال برای اثبات ادعا ثابت می‌کنیم  بعد از رسیدن به $w$ اولین جایی که از مسیر $p$ خارج شویم دیگر به آن برنمیگردیم. 
\\برهان خلف: فرض کنید در جواب بهینه، مطابق شکل زیر بعد از رسیدن به راس $w$، از مسیر $p$ تا راس \متن‌لاتین{s1} برمیگردیم و سپس از مسیر \متن‌لاتین{p} خارج می‌شویم و دوباره در راس \متن‌لاتین{s2} به مسیر \متن‌لاتین{p} برمیگردیم. 
\\
\درج‌تصویر[پهنا=0.4\پهنای‌متن]{A2.PNG}
\\در این صورت هزینه‌ی رفتن از راس \متن‌لاتین{s1} به \متن‌لاتین{s2} در مسیر \متن‌لاتین{p} را \متن‌لاتین{w1} و هزینه‌ی مسیر بین این دو راس که در خارج از \متن‌لاتین{p} طی کرده‌ایم را \متن‌لاتین{w2} می‌نامیم. ادعا می‌کنیم اگر به‌جای مسیری که در خارج از \متن‌لاتین{p} رفته‌ایم از داخل \متن‌لاتین{p} از راس \متن‌لاتین{s1} به راس \متن‌لاتین{s2} می‌رفتیم هزینه‌ی مسیر کمتر می‌شد. چرا که با فرض بهینه بودن مسیری که طی کردیم داریم: 
$w1 \le w2$ 
(در غیر اینصورت در مسیر از \متن‌لاتین{u} به \متن‌لاتین{w} اگر از بخش \متن‌لاتین{w2} می‌رفتیم هزینه کمتر می‌شد که با بهینه بودن مسیر در تناقض است). حال اگر مسیر بین \متن‌لاتین{s1} و \متن‌لاتین{s2} را دوباره در \متن‌لاتین{p} برای دومین بار طی میکردیم، طبق فرض سوال باید هزینه‌ای کمتر از \متن‌لاتین{w1} پرداخت می‌کردیم چون به اندازه‌ی وزن کمتر یال‌ها اینبار هزینه می‌کنیم. از داخل \متن‌لاتین{p} از \متن‌لاتین{s1} به \متن‌لاتین{s2} می‌رفتیم (به جای اینکه از \متن‌لاتین{p}خارج شویم) هزینه‌ای کمتر از \متن‌لاتین{w1} پرداخت می‌کردیم که چون 
$w1 \le w2$
پس هزینه‌ی مسیر جدید کمتر از مسیر قبلی می‌شود و با فرض بهینه بودن مسیری که طی کرده‌ایم در تناقض است. 
بنابراین فرض خلفی که کردیم به تناقض رسید و در نتیجه ادعای ما اثبات می‌شود. 
\\حال برای اینکه مسیر بهینه را بیابیم روی راس \متن‌لاتین{x} (که در شکل اول مشخص شده‌است) حالت‌بندی می‌کنیم و به ازای هر کدام از رئوس گراف بدست می‌آوریم اگر راس \متن‌لاتین{x} ما (یعنی راسی که مسیر ما در آن به سه شاخه منشعب می‌شود) این راس باشد، هزینه‌ی پرداختی چقدر خواهد بود و به ازای تمام حالات کمترین جواب را بدست می‌آوریم. 
\\حال فرض کنید برای راسی مثل \متن‌لاتین{y} می‌خواهیم بدست بیاوریم اگر از \متن‌لاتین{u} به \متن‌لاتین{w} برویم و سپس به راس \متن‌لاتین{y} در مسیر برگردیم و از \متن‌لاتین{y} به \متن‌لاتین{v} برویم هزینه چقدر خواهد بود. در این مسیر یال‌های از راس \متن‌لاتین{u} تا \متن‌لاتین{y} را با وزن بیشتر آن‌ها می‌رویم و بعد برای یال های بین \متن‌لاتین{y} و \متن‌لاتین{w} به اندازه‌ی جمع دو وزن کوچک و بزرگ یال آن‌ها هزینه می‌کنیم و برای یال‌های بین \متن‌لاتین{y} و \متن‌لاتین{v} هم به اندازه‌ی وزن بیشتر آن‌ها هزینه می‌کنیم. بنابر این یک بار روی یال‌های با وزن بیشتر از راس \متن‌لاتین{u} الگوریتم \متن‌لاتین{dijkstra} را اجرا می‌کنیم و به ازای هر راس مثل \متن‌لاتین{i}، کمترین فاصله این راس تا راس \متن‌لاتین{u} را \متن‌لاتین{dis1[i]} می‌نامیم که با \متن‌لاتین{o(nlogn + m)} بدست می‌آید. همچنین یکبار این کار را برای راس \متن‌لاتین{v} انجام می‌دهیم و کمترین فاصله هر راس تا راس \متن‌لاتین{v} را با \متن‌لاتین{dis2[i]} نمایش می‌دهیم (باز هم با \متن‌لاتین{o(nlogn + m)}). یکبار هم وزن هر یال را جمع دو وزن آن قرار داده و روی گراف مذکور از راس \متن‌لاتین{w} الگوریتم \متن‌لاتین{dijkstra} را اجرا می‌کنیم و فاصله‌ها را با \متن‌لاتین{dis3[i]} نمایش می‌دهیم (\متن‌لاتین{o(nlogn + m}).
\\ حال به ازای هر راس \متن‌لاتین{d[i] = dis1[i] + dis2[i] + dis3[i]} را محاسبه می‌کنیم (که \متن‌لاتین{d[i]} یعنی اگر راس وسط ما، که از آن راس مسیر منشعب می‌شود، راس \متن‌لاتین{i} باشد، کمترین هزینه‌ی مسیر چقدر خواهد بود) که با \متن‌لاتین{o(n)} قابل محاسبه است و کمترین \متن‌لاتین{d[i]} جواب مساله است. پس در کل با \متن‌لاتین{o(nlogn + m)} جواب را بدست می‌آوریم.\\
}




%%%%%%%%%%%%%%%%%%%%%%%%%%%%%%%%%%%%%%%%%%%%%%%%%%%%%
% برخی دستورات مفید برای تمرین‌ها. %
%%%%%%%%%%%%%%%%%%%%%%%%%%%%%%%%%%%%%%%%%%%%%%%%%%%%%

% فرمول‌های نمونه:
%$f(x) = \Theta(x)$
%$T(x, y) \leq \log(x) + 2^y + \binom{2n}{n}$
% {\tt code\_function}


% لیست موارد بدون شماره‌گذاری:
%\شروع{فقرات}
%    \فقره مورد اول 
%        \شروع{فقرات}
%            \فقره مورد اول داخلی 
%                \شروع{فقرات}
%                    \فقره اولین مورد خیلی داخلی 
%                    \فقره دومین مورد خیلی داخلی 
%                \پایان{فقرات}
%            \فقره مورد دوم داخلی 
%        \پایان{فقرات}
%    \فقره مورد دوم 
%\پایان{فقرات}

% لیست موارد شماره‌گذاری شده:
%\شروع{شمارش}
%    \فقره یک مورد 
%    \فقره یک مورد دیگر 
%    \فقره باز هم یک مورد دیگر
%\پایان{شمارش}

% برای نوشتن یک فرمول چند خطی که یک شماره بخوره. برای حالت‌های دیگر به سایت زیر مراجعه کنید:
% https://www.overleaf.com/learn/latex/Aligning_equations_with_amsmath
%\begin{align} 
%    \begin{split}
%        (x+y)^3 &= (x+y)^2(x+y) \\
%                &= (x^2+2xy+y^2)(x+y) \\
%                &= (x^3+2x^2y+xy^2) + (x^2y+2xy^2+y^3) \\
%                &= x^3+3x^2y+3xy^2+y^3
%    \end{split}                                 
%\end{align}

% نوشتن ماتریس و آرایه دوبعدی
%\begin{align}
%    A = 
%    \begin{bmatrix}
%        A_{11} & A_{21} \\
%        A_{21} & A_{22}
%    \end{bmatrix}
%\end{align}

% اضافه کردن تصویر با شرح. برای اضافه کردن تصویر بدون شرح، فقط دستور درج‌تصویر را استفاده کنید.
%\شروع{شکل}[h!]
%\تنظیم‌ازوسط
%  \درج‌تصویر[پهنا=0.5\پهنای‌متن]{img.jpg}
%\شرح{شرح تصویر}
%\برچسب{تصویر:۱}
%\پایان{شکل}
%\پایان{زیرمسائل}

%کد زیر را ببینید.
%%\lstinputlisting{source_filename.py}
%\begin{latin}
%\begin{lstlisting}
%# Import the modules
%import sys
%import random
%
%ans = True
%
%while ans:
%question = raw_input("Ask the magic 8 ball a question: (press enter to quit) ")
%
%answers = random.randint(1,8)
%
%if question == "":
%sys.exit()
%
%elif answers == 1:
%print "It is certain"
%
%elif answers == 2:
%print "Outlook good"
%
%elif answers == 3:
%print "You may rely on it"
%
%elif answers == 4:
%print "Ask again later"
%
%elif answers == 5:
%print "Concentrate and ask again"
%
%elif answers == 6:
%print "Reply hazy, try again"
%
%elif answers == 7:
%print "My reply is no"
%
%elif answers == 8:
%print "My sources say no"
%\end{lstlisting}
%\end{latin}